% 第二部分：讨论（由 main_zh.tex 引入）。

\section{讨论}\label{sec:discussion}

\paragraph{共享骨干带来了什么。} 共享骨干上的评判模型与单独加载的评判模型给出相同的答案，而训练器只需保存一份骨干而不是两份：加载后少用 \NUM{mem_saving_load} 显存。对 96\,GB GPU 上的 4B 策略，这只是方便；在更小的 GPU 上或更大的模型上，它可能决定评判模型能否与策略放在一起。代价是每次切换适配器 \NUM{adapter_switch_ms}\,ms，以及评判适配器以未合并方式使用带来的一些速度损失。

\paragraph{评判模型在验证器之上带来了什么。} \texttt{rlvr\_kev\_add} 和 \texttt{rlvr\_kev\_gate} 各有一个种子最终远高于 \texttt{rlvr} 的每个种子，并把各自条件的均值抬到了 \texttt{rlvr} 之上。这两个 run 与其他 run 的区别在于它们改成了先写草稿；本项目中共有 \NUM{alldraftruns} 个 run 这样做了，其中 \NUM{allveronlydraft} 个的奖励只有验证器。同一条件的其他种子最终与 \texttt{rlvr} 相近。每个条件只有 2--3 个种子，我们无法判断奖励中的评判模型是否让这种切换更容易发生（这里加入 \kev{} 分数或用它做门控的 5 个 run 中有 2 个切换，\texttt{rlvr} 的 3 个 run 中没有）。评判模型能起作用的余地本来就很窄。验证器给每个解答一个确切的判定，而 $\lambda<1$ 时评判模型只能在判定相同的解答之间重新排序：在错误解答之间偏向更接近正确的那些，以及在全错的组里——验证器在那里完全不给信号。前几步之后，全错的组已经很少（\texttt{rlvr} 最后 20 步为 \NUM{allwrong_rlvr}），而在判定相同的解答之间，评判模型的偏好可能关乎风格，不亚于关乎正确与否：\kev{} 会为答完之后把题目再解一遍付费（\S\ref{sec:res-repeat}），本项目奖励中含有 \kev{} 分数的 \NUM{allkevruns} 个 run 中有 \NUM{allkevrepeat} 个最终如此，其他 \NUM{allnokevruns} 个 run 中只有 \NUM{allnokevrepeat} 个。这些重复的 run 都没有因此变得更准确。

\paragraph{种子与风格。} 我们的 run 之间最大的差别不在奖励条件之间，而在最终策略的写作风格之间。改成先写草稿的 run 在 AIME24 上达到 \NUM{alldraftaime24min}\% 到 \NUM{alldraftaime24max}\%，本项目其他 run 至多 \NUM{allotheraime24max}\%。切换是突然发生的：在这些 run 的每一个中，MATH-500 解答里草稿所占比例在相隔 50 步的两次评估之间从至多 \NUM{alldraftjumpbeforemax} 升到至少 \NUM{alldraftjumpaftermin}，最早从第 \NUM{alldraftjumpfrommin} 步开始，最晚从第 \NUM{alldraftjumpfrommax} 步开始。验证器对训练解答的判定看不出是什么触发了它。像我们这样每个条件只有一到三个种子的奖励变体比较，可能把这样的切换误当成奖励的效果。因此我们报告每个 run 及其风格；在小模型上比较 RLVR 变体，也应当这样做，或者用足够多的种子来估计每种变体发生切换的频率。先写草稿是更高准确率的原因，还是只是与之相伴，我们没有检验；在评估或训练时强制或禁止草稿开头，会是第一步检验。

\paragraph{评判模型不是验证器。} 离线时，\kev{} 的 $\pc$ 区分基座策略正误解答的 AUROC 接近 0.9（\S\ref{sec:res-study}）。探针表明，不改变答案的改动也能左右它，尽管幅度很小（声称每一步都验证过）。只对 $\pc$ 训练时，策略找到了评判模型高估的解答：它的错误答案得到的 $\pc$ 达到评判模型原先只给正确答案的水平，评判模型在组内的排序退化到接近随机，除 GSM8K 外每个基准上的准确率都跌到基座模型之下。这正是 \citet{gao2023scaling} 所描述的奖励过度优化，只是规模更小、发生在 200 步之内。离线 AUROC 衡量的是评判模型在其训练分布上的表现，而不是在针对它优化的策略所产生的分布上，也不是在策略自己漂移进去的风格上。

\paragraph{局限。} 一种策略与评判模型规模（4B），一个任务（MATH，在 MATH-500、GSM8K 和三个竞赛题集上评估），200 步，一个为 \texttt{rlvr} 选定的学习率，一个评判权重，每个条件 2--3 个种子——不足以说明一个条件以多大频率落到每种风格。区间按题目重采样，不按种子。vLLM 上对同一模型的重复贪心评估相差约 \NUM{szsd_math500} 个百分点（\S\ref{sec:res-cond}）。两个风格量是在读过最终解答之后定义的。仓库中没有策略的权重（有它们的解答、日志和训练解答），因此无法从仓库出发对最终策略做干预实验。训练期间其他实验共用 GPU，所以除在空闲 GPU 上测得的时间外，我们不报告训练时间。

\paragraph{与早期版本的关系。} 本文第 1--3 版报告过一个排行榜，其中 \kev{} 的奖励优于验证器，还报告过由假设系数算出的显存节省。两者都没有被测量过（见文首声明和 \file{CORRECTION.md}）。实测的答案是：共享骨干省下一份骨干，且不改变评判模型；在我们的 run 中，把 \kev{} 的判断加入验证器并没有稳定地改进 MATH 上的 GRPO。它的两个 run 确实大幅提高了，但伴随着一种写作风格的改变，而没有评判模型的 run 也发生过这种改变。
